%----------SUMMARY----------%
\section{요약}
\begin{itemize}[leftmargin=0.15in, label={}]
    \small{\item{
    온프레미스 Kubernetes 위에 MLOps 플랫폼을 직접 설계·구축·운영하는 엔지니어입니다. HUINNO에서 6개월간 수집 → 레이크하우스 → 변환 → 계보 → 분산 학습 → 실험·모델 레지스트리까지 전 구간을 \textbf{Pulumi}로 코드화했고 플랫폼 2인 팀에서 \textbf{클러스터 구성은 단독}으로 맡았습니다(인프라 스택 코드 변경의 93\%, ETL 저장소 커밋의 83\%를 직접 작성). 도구를 늘리기보다 사람이 병목인 자리를 찾아 없애는 쪽을 택하고, AI 코딩 에이전트를 붙여 한 사람이 내놓는 양을 팀 규모로 끌어올립니다. 그 전에는 iOS 개발자로 시작해 백엔드·프론트·클라우드 데이터 파이프라인을 거쳐 데이터 엔지니어링 팀을 리드했습니다.
    }}
\end{itemize}

%----------KEY RESULTS----------%
\section{핵심 성과 (2026-09 실측)}
\begin{itemize}[leftmargin=0.15in, label={}]
    \small{
    \item{\textbf{MLOps 파이프라인 전체 설계와 구현을 주도}. 수집부터 모델 레지스트리까지 끊긴 데 없이 이어 붙이고 12개 Pulumi 스택으로 코드화 — 인프라를 고쳐도 PR로 남고 새 환경은 스택만 골라 다시 세울 수 있음. (저장소 커밋 1,204건 · 머지 PR 678건)}
    \item{\textbf{10노드 · 상시 pod 235개 · 사내 서비스 23종}을 프로덕션 운영. 23종 전부에 Dex(GitHub OIDC) + oauth2-proxy SSO를 적용하고 같은 신원을 Kubernetes RBAC까지 연결.}
    \item{\textbf{데이터셋 준비를 한 사람에게 기대지 않게} — 담당자에게 부탁하고 기다리던 일을 Airflow 3.x + dbt + Apache Iceberg가 정해진 때에 알아서 돌리도록 바꿈. DataHub에 출처를 자동으로 남겨 연구자가 직접 찾아 쓰게 함. (ETL 저장소 커밋 1,599건 · 공용 프레임워크와 제품 파이프라인 · MIMIC-IV · 운영 파이프라인을 직접 작성)}
    \item{\textbf{팀 하나가 붙을 일을 6개월에}. Kubernetes MLOps 플랫폼을 AI 코딩 에이전트와 함께 대부분 직접 짜 지금도 프로덕션으로 돌리고 있음. 속도만 낸 게 아니라 SDK 단위 테스트 719건 · CI 게이트 · 실패 기록 23건 · \mbox{PRD 30건}을 같이 쌓음.}
    \item{\textbf{이전 경력:} \mbox{AWS EMR/PySpark} 미술 시장 \mbox{ETL 파이프라인} 총괄(Eazel, 팀 리드).}
    }
\end{itemize}

%----------CURRENT WORK----------%
\section{현재 하는 일}
\begin{itemize}[leftmargin=0.15in, label={}]
    \small{\item{
    HUINNO Data Engineer (2025.12 --). \textbf{연구가 빨라지도록 받치는 플랫폼을 설계·구축하고, 그 위에 MLOps 파이프라인을 하나씩 얹고 있습니다.} \\[3pt]
    베어메탈 Kubernetes 위에서 네트워크 · 스토리지 · 카탈로그 · 계보 · 오케스트레이션 · 분산 학습 · 신원까지 전 층을 운영합니다. \\[3pt]
    지금은 연구 코드가 KubeRay 분산 학습으로 넘어오도록 재현성 계약과 MLflow 트래킹 · 모델 레지스트리를 붙이는 중입니다. 이 계약은 데이터셋에 버전 · 내용 해시 · 코드 커밋 · 환경 lock · 시드를 실행 시점에 박아 둡니다. \\[3pt]
    기존 연구 워크플로우를 정량 조사해 보니 결과만 남고 그걸 만든 데이터 · 코드 · 환경은 어디에도 붙어 있지 않았습니다. 거기서 이 설계가 출발했습니다. \\[3pt]
    「도메인 ETL은 연구자가, 실행 · 계보 · 검증 · 재시도는 플랫폼이」 맡는 분담을 정착시키고 있습니다.
    }}
\end{itemize}

%----------STACK----------%
\section{기술 스택}
\begin{itemize}[leftmargin=0.15in, label={}]
    \small{\item{
    \textbf{플랫폼}: Kubernetes(kubeadm) · Pulumi · Cilium · SeaweedFS · Dex/OIDC · Prometheus/Grafana \\[2pt]
    \textbf{데이터}: Airflow 3.x · dbt · Apache Iceberg · DataHub · PostgreSQL(pg\_duckdb) · ClickHouse \\[2pt]
    \textbf{ML}: Ray/KubeRay · MLflow · PyTorch · DVC \\[2pt]
    \textbf{언어·기타}: Python · SQL · Swift · Shell · AWS/GCP · LangChain
    }}
\end{itemize}

%----------CERTIFICATIONS----------%
\section{자격증 · 어학}
\begin{itemize}[leftmargin=0.15in, label={}]
    \small{\item{
    \textbf{자격증}: 정보처리기사 \mbox{(2017.11)} · 정보처리운용기능사 \mbox{(2009.07)} — 한국산업인력공단 발급 \\[2pt]
    \textbf{어학}: 한국어 원어민 · 영어 \mbox{OPIc AL} · 일본어 일상 회화
    }}
\end{itemize}

%----------DETAIL----------%
% 섹션 제목을 붙이기엔 한 문장뿐이라 각주 한 줄로 둔다
\vspace{6pt}
{\small 경력 · 프로젝트 상세는 별도 문서(\texttt{resume\_korean.pdf})에 있습니다.}
